\subsection{R-23: the exact homological defect after restriction to four variables}

\paragraph{Statement R-23.}
Let $m\geq3$, let $R=\mathbb C[X_{ij}:1\leq i,j\leq m]$, and let
$J_m$ be the ideal generated by all $(m-1)\times(m-1)$
subpermanents, equivalently the first partial derivatives of
$\operatorname{per}_m$. Let
\[
 S=R/(L_1,\ldots,L_{m^2-4})\simeq\mathbb C[x,y,z,w]
\]
be a general four-variable linear restriction. Write
\[
 \epsilon_m=\dim_{\mathbb C}
       \operatorname{Tor}_1^R(R/J_m,S)_m.
\]
Then
\begin{equation}\label{R23:Tor_defect}
 \dim\mathcal P_{m,4}=4m^2-2m+2-\epsilon_m.
\end{equation}
Moreover
\begin{equation}\label{R23:first_degree}
 \operatorname{Tor}_1^R(R/J_m,S)_d=0\qquad(d<m),
\end{equation}
and R-22 implies
\[
 0\leq\epsilon_m\leq20m^{3/2}-2m\qquad(m\geq99).
\]
Thus attaining the exact row--column torus bound is equivalent to one
specified degree of one Tor group vanishing. This vanishing is not
proved here for $m\geq17$.

The general-section defect is already nonzero in a complete known
range. By R-13, for $3\leq m\leq16$ it is exactly
\begin{equation}\label{R23:small_defect}
 \epsilon_m=4m^2-2m+2-\binom{m+3}{3}>0.
\end{equation}
Consequently the full-variable linear syzygy theorem cannot simply be
carried through a general four-variable restriction. Nor can this
restriction be justified by treating all its defining linear forms as
a regular sequence on $R/J_m$: this is impossible for every $m\geq3$.

\begin{proof}
We first specify the cited full-variable input. Efremenko, Landsberg,
Schenck, and Weyman (2018), Theorem~1.1 and \S4.1, Example~4.1
and Theorem~4.3 in the author version, determine the linear strand
of the subpermanent ideal (R-B16). At subpermanent size $m-1$, the
space of linear first syzygies has dimension $2m-2$. A basis consists
of differences between row Laplace expansions of the full permanent,
and differences between column expansions. In the minimal resolution
of $R/J_m$, therefore,
\begin{equation}\label{R23:resolution_start}
 F_1=R(-(m-1))^{m^2},\qquad
 (F_2)_m=\mathbb C^{2m-2},
\end{equation}
and $F_2$ has no components below degree $m$. Other generators of
$F_2$ have higher degree and do not enter the calculation below.

\paragraph{The old linear relations remain independent.}
Write $G_{ij}=\partial\operatorname{per}_m/\partial X_{ij}$.
The row and column expansion vectors are
\[
 \mathsf R_i=(\delta_{ui}X_{uv})_{u,v},\qquad
 \mathsf C_j=(\delta_{vj}X_{uv})_{u,v}.
\]
Each maps to the permanent under $(f_{ij})\mapsto\sum f_{ij}G_{ij}$.
Thus the $2m-2$ differences
$\mathsf R_i-\mathsf R_1$ and $\mathsf C_j-\mathsf C_1$ are
the cited syzygies. A linear combination is the vector
$((\alpha_i+\beta_j)X_{ij})$, where
$\sum\alpha_i=\sum\beta_j=0$.

After restriction to the cyclic pencil of R-15, both its diagonal
and its cyclic superdiagonal entries are nonzero linear forms.
If this relation vector vanishes, its diagonal entries give
$\alpha_i+\beta_i=0$, and its cyclic entries give
$\alpha_i+\beta_{i+1}=0$. Hence all $\beta_i$ are equal and all
$\alpha_i$ are their negatives. Their zero sums force every
coefficient to vanish. The old relations are therefore independent
on that pencil, and hence on a nonempty Zariski open set of
four-variable restrictions.

\paragraph{There are no new constant relations generically.}
For completeness we prove independence of the $m^2$ restricted
gradient generators on the same cyclic pencil. Number indices
$0,\ldots,m-1$, take a primitive $m$th root $q$, and put
$\ell_i=x-q^iy$, $v_i=z-q^iw$. Each complementary permanent has
one matching. For each $0\leq k\leq m-1$, its $m$ complementary
permanents of $(z,w)$ degree $k$ are the cyclic translates of
\begin{equation}\label{R23:cofactor_orbit}
 F_k=\prod_{i=0}^{m-k-2}\ell_i
                 \prod_{i=m-k-1}^{m-2}v_i.
\end{equation}
For example this is the cofactor deleting row $m-1$ and column
$m-k-1$. Empty products have value one.

Every coefficient of each of the two geometric binary products in
\eqref{R23:cofactor_orbit} is nonzero. Indeed each length is below
$m$, so the geometric binomial-product formula has no factor
$1-q^j$ with $j$ divisible by $m$. For every
$0\leq t\leq m-1$, choose integers
$0\leq a\leq m-k-1$, $0\leq b\leq k$ with $a+b=t$.
The corresponding coefficient of $y^aw^b$ is nonzero. Under
$\tau:(y,w)\mapsto(qy,qw)$ this gives a nonzero component of
character $t$. Thus the orbit
$F_k,\tau F_k,\ldots,\tau^{m-1}F_k$ has dimension $m$ by Fourier
inversion. The different $k$ have distinct bidegrees, proving
independence of all $m^2$ cofactors. Their restrictions are therefore
independent for general $S$.

\paragraph{The exact Tor identity.}
Tensor a minimal free resolution of $R/J_m$ with $S$ and take
homogeneous degree $m$. The part relevant to its first homology is
\begin{equation}\label{R23:degree_complex}
 \mathbb C^{2m-2}\longrightarrow
                 S_1^{m^2}\longrightarrow S_m.
\end{equation}
The first map is injective on the general open set just proved.
The second sends $(f_{ij})$ to
$\sum f_{ij}G_{ij}|_S$. It is exactly the differential of the
permanent projection map, because varying any pencil entry by a
linear form $f_{ij}$ changes the permanent by that expression.
In characteristic zero its rank at a general pencil equals the
dimension $c_m$ of the image closure. Consequently the first
homology of \eqref{R23:degree_complex} has dimension
$4m^2-c_m-(2m-2)$, proving \eqref{R23:Tor_defect}.
In degrees below $m-1$ the module $F_1\otimes S$ is zero. In degree
$m-1$ its map to $S$ is injective by the cofactor independence just
proved. This gives \eqref{R23:first_degree}.
Substitution of R-22 proves the stated upper bound for $\epsilon_m$.

\paragraph{Why the known linear strand does not supply the missing vanishing.}
R-13 gives $c_m=\binom{m+3}{3}$ for $3\leq m\leq16$,
and hence \eqref{R23:small_defect}. The positive values are, in order,
\[
 12,\ 23,\ 36,\ 50,\ 64,\ 77,\ 88,\ 96,\ 100,\ 99,\ 92,\ 78,\ 56,\ 25.
\]
These are general-section Tor groups in the very same permanent
family, even though the complete full-variable linear syzygy space
has only the $2m-2$ displayed generators. The larger-section
vanishing sought for $m\geq17$ needs an additional theorem.
At the special cyclic section the defect is even larger:
\[
 \dim\operatorname{Tor}_1^R(R/J_m,S_{\mathrm{cyc}})_m
 =2m(m-1),
\]
because its differential image has dimension $2m^2+2$ and the old
relations remain independent. This special value is not substituted
for the general-section defect.

Finally Boralevi, Carlini, Micha{\l}ek, and Ventura (2024),
Theorem~4.14, record von zur Gathen's bound
$\operatorname{ht}J_m\geq5$ for every $m\geq3$; their Theorem~4.23
strengthens it to six when $m\geq6$ (R-B17). Therefore
\[
 \operatorname{depth}(R/J_m)
 \leq\dim(R/J_m)\leq m^2-5<m^2-4.
\]
A regular sequence on $R/J_m$ cannot have the required length
$m^2-4$. Thus tensoring its resolution with the four-variable
quotient is not justified as an exact operation by a regular-sequence
argument. This obstruction does not itself decide the single graded
piece in \eqref{R23:Tor_defect}: higher-degree Tor may survive even
if that piece vanishes.
\end{proof}

